%! The Quadratic Formula & the Discriminant — Unit 3, Lesson 3.5 — Lesson Plan
\documentclass[10pt]{article}
\usepackage{algebra2-article}
\usepackage{algebra2-boxes}
\usepackage{graphicx}   % -article does NOT load graphicx; needed for thumbnails/screenshots
\graphicspath{{images/}}
\newcommand{\TallMath}[1]{$\displaystyle #1\rule[-1.4em]{0pt}{3.2em}$}
\newcommand{\CourseName}{Algebra 2: Shepherd}
\newcommand{\SchoolYear}{2026--2027}
\newcommand{\MeetingLength}{55 minutes}
\newcommand{\UnitNumberName}{Unit 3: Quadratic Functions \quad}
\newcommand{\LessonNumberName}{Lesson 3.5: The Quadratic Formula \& the Discriminant}

\begin{document}
\begin{center}
    {\Large\bfseries \CourseName: \SchoolYear \\
    {\small\bfseries \UnitNumberName \LessonNumberName}}
\end{center}

\begin{tcolorbox}[colback=forestbg, colframe=forest, boxrule=0.9pt, arc=2mm,
    left=2mm, right=2mm, top=1.5mm, bottom=1.5mm]
    \textbf{Primary Objective:} Students learn the \emph{one method that solves every quadratic}. They
    derive the \textbf{quadratic formula} $x=\dfrac{-b\pm\sqrt{b^2-4ac}}{2a}$ by completing the square on
    the general equation $ax^2+bx+c=0$ (the Lesson~3.3 move, done once in general), then use it fluently:
    write the equation in standard form, read off $a,b,c$ (with signs), and simplify --- including
    \emph{complex} answers when the radicand is negative (Lesson~3.4). They isolate the \textbf{discriminant}
    $b^2-4ac$ as a \emph{predictor}: its sign tells them, \emph{before} solving, whether a quadratic has two
    real roots ($>0$), one repeated real root ($=0$), or two complex-conjugate roots ($<0$) --- and they
    connect each case to the parabola's $x$-intercepts (two / one / none). They choose a solving method
    strategically, verify solutions by substitution, and apply the formula to model a contextual situation.
    \\[2pt]
    \textbf{Standards (2023 VA SOL):} \textbf{A2.EI.2b} (solve a quadratic equation in one variable over the
    set of complex numbers algebraically --- the quadratic formula), \textbf{A2.EI.2a} (create a quadratic
    equation to model a contextual situation), and \textbf{A2.EI.2d} (verify solutions algebraically and
    interpret them in context). The discriminant's link to the number and type of $x$-intercepts revisits
    \textbf{A2.F.2d}. Builds directly on completing the square (\textbf{A2.EI.2b}, Lesson~3.3) and the
    imaginary unit $i$ (\textbf{A2.EO.4}, Lesson~3.4).
\end{tcolorbox}

\renewcommand{\arraystretch}{1.4}
\begin{skillbox}[Priority Ideas \& Skills]{goldbox}
\begin{tabularx}{\linewidth}{@{}X|X@{}}
\textbf{Priority Ideas \& Skills} & \textbf{Key Understandings (the ``why'')} \\[2pt]
\hline
\vspace{-6pt}
\begin{itemize}[leftmargin=1.1em, nosep, after=\vspace{2pt}]
  \item Write $ax^2+bx+c=0$ in standard form and identify $a,b,c$ with correct signs.
  \item Apply $x=\dfrac{-b\pm\sqrt{b^2-4ac}}{2a}$ and simplify to simplest radical or $a+bi$ form.
  \item Compute the \textbf{discriminant} $b^2-4ac$ and classify the roots \emph{before} solving.
  \item Match each discriminant case to the parabola's number of $x$-intercepts (2 / 1 / 0).
  \item Choose a solving method strategically; verify a solution by substitution.
  \item Model a contextual situation with a quadratic and solve it with the formula.
\end{itemize}
&
\vspace{-6pt}
\begin{itemize}[leftmargin=1.1em, nosep, after=\vspace{2pt}]
  \item The formula is completing the square done \emph{once, in general} --- so it works for \emph{every} quadratic.
  \item A negative radicand is no longer a dead end: Lesson~3.4's $i$ turns it into a complex-conjugate pair.
  \item One number, $b^2-4ac$, encodes the entire ``how many / what kind'' of the roots.
  \item Roots \emph{are} $x$-intercepts, so ``two complex roots'' means the parabola misses the $x$-axis.
  \item Factoring is fastest \emph{when it works}; the formula is the universal fallback.
  \item Not every algebraic solution is physically meaningful --- interpret and reject impossible answers.
\end{itemize}
\end{tabularx}
\end{skillbox}

\begin{skillbox}[{Vocabulary, Concepts, \& Theorems}]{sky}
\renewcommand{\arraystretch}{1.3}
\begin{tabularx}{\linewidth}{@{}l X@{}}
\textbf{Quadratic formula} & $x=\dfrac{-b\pm\sqrt{b^2-4ac}}{2a}$; solves \emph{any} $ax^2+bx+c=0$ ($a\ne0$) over the complex numbers. \\
\textbf{Standard form} & $ax^2+bx+c=0$ --- everything on one side, $=0$ --- required before reading $a,b,c$. \\
\textbf{Discriminant} & $b^2-4ac$, the expression under the radical; its sign predicts the nature of the roots. \\
\textbf{Two real roots} & $b^2-4ac>0$: two distinct real solutions; parabola crosses the $x$-axis twice. \\
\textbf{Double (repeated) root} & $b^2-4ac=0$: one real solution counted twice; vertex sits \emph{on} the $x$-axis. \\
\textbf{Complex-conjugate roots} & $b^2-4ac<0$: two solutions $a\pm bi$; parabola never touches the $x$-axis. \\
\end{tabularx}
\end{skillbox}

\begin{fixedskillbox}[Activate Prior Knowledge \& Spiral Review]{forestbg}
    The Warm-Up rehearses the three moves today's formula requires: \textbf{(1)} \emph{evaluating}
    $b^2-4ac$ for given $a,b,c$ (signed arithmetic --- the single most error-prone step, and today it
    \emph{is} the discriminant); \textbf{(2)} \emph{simplifying radicals}, including negative ones from
    Lesson~3.4 ($\sqrt{-16}=4i$, $\sqrt{20}=2\sqrt5$), so formula outputs land in simplest radical or
    $a+bi$ form; and \textbf{(3)} \emph{rewriting an equation in standard form} $ax^2+bx+c=0$ and reading
    off $a,b,c$ with signs. Debrief item~1 last and connect it forward: ``the number you just computed has
    a name, and it predicts the answer before you solve.''
\end{fixedskillbox}

\begin{skillbox}[Hook]{forestbg}
    Put four quadratics on the board and challenge students to solve them: $x^2-2x-3=0$ (factors),
    $x^2-6x+7=0$ (won't factor --- needs completing the square), $x^2+2x+5=0$ (complex roots, from 3.4),
    and $2x^2-7x+3=0$. Ask: ``Which \emph{method} for each?'' Let the friction build --- factor, square
    roots, complete the square, each with its own conditions. Then reframe: ``What if \emph{one} formula
    solved all four, every time, no guessing?'' Reveal that completing the square on the \emph{general}
    equation $ax^2+bx+c=0$ produces exactly that formula --- and, as a bonus, one piece of it,
    $b^2-4ac$, tells you what kind of answer you'll get \emph{before} you turn the crank.
\end{skillbox}

\begin{skillbox}[Lesson]{forestbg}
\begin{multicols}{2}
\textbf{1. Derive the formula (once, in general).} Complete the square on $ax^2+bx+c=0$ --- exactly the
Lesson~3.3 move, but with letters. Land the result and box it:
\[ x=\frac{-b\pm\sqrt{b^2-4ac}}{2a}. \]
Emphasize: because it came from completing the square, and \emph{that} works on any quadratic, this
formula works on \emph{every} quadratic. No more choosing whether a method applies.

\textbf{2. Use it.} The protocol: (i) write $ax^2+bx+c=0$ in \emph{standard form}; (ii) identify $a,b,c$
\emph{with signs}; (iii) substitute and simplify. Work the unit anchor $x^2-2x-3=0$: $a=1,b=-2,c=-3$,
so $x=\dfrac{2\pm\sqrt{16}}{2}=\dfrac{2\pm4}{2}=3,-1$ --- the same zeros as factoring $(x+1)(x-3)$ in 3.2
and reading the graph in 3.0. \emph{Same curve, same roots, universal method.}

\columnbreak

\textbf{3. The discriminant.} Isolate $D=b^2-4ac$ and make it a \emph{predictor}. Build the three-case
table with a pre-drawn parabola for each: $D>0\Rightarrow$ two real roots (two $x$-intercepts);
$D=0\Rightarrow$ one repeated real root (vertex on the axis, one $x$-intercept); $D<0\Rightarrow$ two
complex-conjugate roots (no $x$-intercepts). Stress: you can answer ``how many / what kind'' \emph{without}
finishing the solve.

\textbf{4. Complex solutions.} Run $x^2+2x+5=0$: $D=4-20=-16<0$, so
$x=\dfrac{-2\pm\sqrt{-16}}{2}=\dfrac{-2\pm4i}{2}=-1\pm2i$ --- the \emph{same} conjugate pair the
completing-the-square capstone produced in 3.4, now from the formula. Reinforce: a negative discriminant
is not ``no solution,'' it is a \emph{complex} solution.

\textbf{5. Choose a method \& model.} Factoring is fastest when it works; square roots for $x^2=k$; the
formula is the universal fallback. Close with a projectile/area model (A2.EI.2a/d): create the equation,
solve with the formula, reject the impossible root, interpret in context.
\end{multicols}
\end{skillbox}

\begin{skillbox}[Explicit Instruction: Solving with the Quadratic Formula]{forestbg}
\begin{tabularx}{\linewidth}{@{}p{0.46\linewidth} X@{}}
\textbf{Steps (post and refer back):}
\begin{enumerate}[leftmargin=1.4em, nosep, itemsep=2pt]
  \item Write it in standard form $ax^2+bx+c=0$ (everything on one side).
  \item Identify $a,b,c$ \emph{with signs}.
  \item Compute the discriminant $b^2-4ac$ first --- it previews the answer.
  \item Substitute into $x=\dfrac{-b\pm\sqrt{b^2-4ac}}{2a}$ and simplify (radical or $a+bi$).
\end{enumerate}
&
\textbf{Worked example.} Solve $3x^2=2x+4$.
\begin{itemize}[leftmargin=1.2em, nosep, itemsep=1pt]
  \item Standard form: $3x^2-2x-4=0$, so $a=3,\ b=-2,\ c=-4$.
  \item Discriminant: $(-2)^2-4(3)(-4)=4+48=52>0$ (two real).
  \item $x=\dfrac{2\pm\sqrt{52}}{6}=\dfrac{2\pm2\sqrt{13}}{6}$.
  \item Simplify: $x=\dfrac{1\pm\sqrt{13}}{3}$.
\end{itemize}
\end{tabularx}
\end{skillbox}

\begin{skillbox}[Active Monitoring]{redbox}
Circulate during the notes practice and Tier work. Watch for and correct:
\begin{itemize}[leftmargin=1.2em, nosep]
  \item skipping standard form --- reading $a,b,c$ before the equation is set $=0$;
  \item sign errors in $b^2-4ac$, especially $-4ac$ when $c<0$ (the ``$+$'' the double negative creates);
  \item forgetting $-b$ means the \emph{opposite} of $b$ (so $b=-2$ gives $-b=+2$);
  \item dividing only \emph{part} of the numerator by $2a$ (must divide \emph{both} terms / simplify the whole fraction);
  \item stopping at a negative radicand instead of converting with $i$ (Lesson~3.4);
  \item leaving radicals or fractions unsimplified ($\sqrt{52}=2\sqrt{13}$; reduce common factors).
\end{itemize}
Cold-call prompts: ``What must be true before you read $a,b,c$?'' ``$b=-5$ --- what is $-b$?''
``The discriminant is $-16$; how many real roots, and what kind are they?'' ``Where did this formula
\emph{come from}?''
\end{skillbox}

\begin{skillbox}[Group Work \& Differentiation]{redbox}
\textbf{One Formula, Every Quadratic} activity (tiered; mirrors the notes).
\begin{multicols}{3}
\textbf{Tier R --- Remediate.} Put equations in standard form, identify $a,b,c$, compute the discriminant,
and classify the roots (two real / one repeated / two complex) --- \emph{without} finishing the solve ---
matching each to a pre-drawn parabola.

\columnbreak
\textbf{Tier A --- Approaching.} Solve with the formula: real rational, real irrational (simplest radical
form), and complex-conjugate cases; verify one root by substitution.

\columnbreak
\textbf{Tier E --- Extension.} Model a contextual situation (projectile/area), reject the impossible root
and interpret; and a discriminant \emph{parameter} problem --- find the $k$ that makes an equation have
exactly one real root.
\end{multicols}
\end{skillbox}

\begin{skillbox}[Individual Work \& Assessment]{redbox}
\textbf{Exit Ticket} (independent, no notes): solve $x^2-4x+13=0$ with the formula (complex answer
$2\pm3i$); use the discriminant to state the nature of the roots of a second equation \emph{without}
solving; and an \textbf{SOL-style multiple-choice} item reading the root type from a given discriminant.
Collect and sort into ``got it / sign slip in $b^2-4ac$ / dropped the $\pm$ or the $i$ / didn't simplify''
to decide whether Lesson~3.6 opens with a re-teach.
\end{skillbox}

\begin{skillbox}[Reinforcement \& Extension]{goldbox}
\textbf{Homework:} identify $a,b,c$; compute discriminants and classify roots; solve with the formula
(real rational, real irrational, and complex cases); match discriminant sign to a graph's $x$-intercepts;
choose an efficient method; and verify a solution by substitution. \textbf{Extension:} a projectile model
solved with the formula and a find-the-parameter discriminant problem. \textbf{Preview:} Lesson~3.6,
\emph{Systems Involving Quadratics} --- solving linear--quadratic and quadratic--quadratic systems, where
substitution lands you back on a single quadratic to solve (often with today's formula).
\end{skillbox}
\end{document}
